\documentclass{article}

% Paquetes
\usepackage[utf8]{inputenc} % Para caracteres en español
\usepackage[spanish]{babel} % Para la correcta separación de sílabas y etiquetas en español
\usepackage{graphicx}
\usepackage{amssymb}

\title{}
\author{}
\date{}

\begin{document}

\maketitle

\section{Análisis de las tecnologías Frontend}
En la actualidad existen multitud de opciones en el mercado para el desarrollo de tecnologías móviles. A continuación se describen las tecnologías más populares y utilizadas en la actualidad para el desarrollo de aplicaciones móviles.

\subsection{Desarrollo nativo}
En primer lugar tenemos las tecnologías de \textbf{desarrollo nativo}, es decir, las tecnologías que permiten la creación de aplicaciones específicamente diseñadas
para un sistema operativo en concreto. Dichas tecnologías son: 

\begin{itemize}
    \item \textbf{Kotlin.} Kotlin es un lenguaje que permite desarrollar aplicaciones para todas las plataforams soportadas por Android, siendo Android Studio
    el entorno de desarrollo integrado (IDE) principal para la creación de aplicaciones nativas en dicho sistema operativo. Una de las ventajas más destacas es su 
    total interoperabilidad con Java, lo que facilita el uso de bibliotecas escritas en Java dentro de proyectos desarrollados en Kotlin. Además, el lenguaje proporciona mecanismos
    avanzados de seguridad frente a valores nulos, lo que ayuda a prevenir errores durante la ejecución de las aplicaciones. Entre sus deventajas destaca que, aunque la popularidad ha crecido, aún no alcanzar
    el nivel de adopción de Java, lo que puede dificultar encontrar soluciones a problemas específicos. Asimismo, aunque existen iniciativas para su uso multiplataforma, Kotlin sigue estando principal,emte 
    orientado al desarrollo de aplicaciones en Android.
    \item \textbf{Swift.} Es un lenguaje, diseñado específicamente para crear aplicaciones de alto rendimiento y con una sintaxis clara y moderna. En este caso el IDE es XCode utilizado para el desarrollo nativo en las plataformas de Aplle, como iOS y macOS. 
    Swift se caracteriza por su seguridad, facilidad de uso y capacidades de optimización, lo que lo convierte en una buena opción para el desarrolo de aplicaciones en el ecosistema de Apple. Como principales carácteristicas de Swift podemos destacar una \textbf{sintaxis simple y moderna}, 
    un \textbf{alto rendimiento} debido a su compilador avanzado permite ejecutar aplicaciones de manera eficiente, \textbf{seguridad} gracias a su sistema de tipos fuertes y a la prevención de errores en tiempo de ejecución, e \textbf{interoperabilidad con Objetive-C} lo que facilita la integración de código. 
    Sin embargo, una de las principales deventajas de Swift es similar a la mencionada anteriormente en Kotlin, \textbf{la duplicidad de esfuerzos} en el desarrollo de aplicaciones multiplataforma. Las aplicaciones desarrolladas con Swift
    obligan a los desarrolladors a crear aplicaciones independientes para otras plataformas. 
\end{itemize}

\subsection{Desarrollo Híbrido}
Las tecnologías de desarrollo híbrido representan una alternativa eficaz en el ámbito de aplicaciones multiplataforma. Estas tecnologías permiten crear una única 
base de código que puede ser ejecutada tanto en dispositivos con sistema operativo iOS como en Android. La principal ventaja de este enfoque es la optimización de los recursos,
tanto económicos como operativos, ya que reduce significativamente los costes asociados al desarrallo y mantenimiento de las aplicaciones. Al utilizar un solo lenguaje y una 
única estructura de programación, se consigue implementar una solución funcional para ambas plataformas. Además, permite una mejor experiencia al usuario, garantizando que la aplicación
tenga un comportamiento y apariencia similares. Algunos ejemplos de esta tecnología son:

\begin{itemize}
    \item \textbf{Ionic.} Es un framework de código abierto que permite crear aplicaciones móviles, web y de escritorio utilizando tecnologías webs, como HTML, CSS y JavaScript.
    Se integra con diferentes frameworks frontend como Angular React o Vue, lo que permite a los desarrolladores utilizar el conjunto de herramientas y librerías de su preferencia. Además, Ionic 
    tiene componentes nativos que proporcionan una biblioteca de elementos visuales optimizados para dispositivos móviles. Cuenta también con una \textbf{amplia comunidad} y herramientas adicionales para la integración 
    con servicios de backend y otras tecnologías.
    \item \textbf{Flutter.} Flutter es un framework de código abierto desarrollado por Google que permite crear aplicaciones nativas multiplatadorma con una única base de código. Utiliza el lengujae de programación Darrt
    y proporciona un enfoque único para desarrollar aplicaciones móviles, web y de escritorio. Tiene un \textbf{alto rendimiento} debido a que utiliza un motor gráfico y compila directamente a código nativo, lo que garantizando
    un rendimiento similar al de aplicaciones nativas. Además. cuenta con \textit{Hot Reload} que facilita el desarrollo rápido al permitir que lo cambios en el código se reflejen instantáneamente en la aplicación
    sin tener que recompilar por completo. 
    \item \textbf{React Native.} React Native es un framework de código abierto desarrollado por Facebook que permite crear aplicaciones móviles nativas utilizando JavaScript y React. A diferencia de otros frameworks
    híbridos, React Native no se limita a emular una experiencia nativa, en su lugar, \textbf{genera componentes nativos}, lo que permite un rendimiento más cercano al de una aplicación pura. Además cuenta con una 
    \textbf{amplia comunidad} y una gran cantidad de librerías y herramientas que facilitan el desarrollo de aplicaciones.
    \item \textbf{Xamarin.} Es un framework de código abierto desarrollado por Microsoft que permite crear aplicaciones móviles nativas para iOS y Android utilizando C\# y .NET. Xamarin permite compartir hasta un 90\% del código.
    Además, proporciona acceso a las API nativas de cada plataforma, lo que permite a los desarrolladores crear aplicaciones con una \textbf{experiencia de usuario nativa}. También se integra a la perfección con herramientas del ecosistema de Microsoft, como Visual Studio, facilitando
    el desarrollo y la depuración de aplicaciones.

\end{itemize}



\subsection{Tabla comparativa de tecnologías}

\begin{table}[ht]
    \centering
    \scalebox{0.9}{
    \begin{tabular}{|c|c|c|c|c|c|}
        \hline
        \textbf{Tecnología} & \textbf{Lenguaje} & \textbf{IDE} & \textbf{Ventajas} & \textbf{Desventajas} & \textbf{Popularidad} \\
        \hline
        Kotlin & Java & Android Studio & Interoperabilidad con Java & Baja adopción & Media \\
        \hline
        Swift & Swift & XCode & Sintaxis simple y moderna & Duplicidad de esfuerzos & Alta \\
        \hline
    \end{tabular}
    }
    \caption{Tabla comparativa de tecnologías}
    \label{tab:comparativa}
\end{table}



\end{document}

